% EDITORIAL QUESTION: Como correlacionar de forma unificada as escalas internacionais e estruturar uma matriz cromática de nove faixas protegendo a iniciação?
\section{Matrizes de Equivalência Internacional e Referência de Circuitos}

Esta seção apresenta a \textbf{Tabela~\ref{tab:equivalencia_mestra}} (Matriz Mestra de Equivalência Internacional com foco estrito em boulder e lance-crux) e a \textbf{Tabela~\ref{tab:circuitos_referencia}} (Matriz de Referência de Nove Faixas Cromáticas).

\begin{table}[!htbp]
\centering
\caption{Matriz Mestra de Equivalência Internacional (Foco em Boulder e Lance-Crux).}
\label{tab:equivalencia_mestra}
\footnotesize
\renewcommand{\arraystretch}{1.15}
\begin{tabularx}{\linewidth}{ccccccX}
\toprule
\textbf{IRCRA} & \textbf{V-Scale} & \textbf{FB} & \textbf{CBME} & \textbf{YDS} & \textbf{Ewbank} & \textbf{Exigência Biomecânica e Perfil Motor Típico} \\
\midrule
\textbf{8}  & —   & 4-  & IV    & 5.8   & 14 & Iniciação absoluta; subida tipo escada em jarros (\textit{jugs}) amplos e simétricos. \\
\rowcolor{rowgray}
\textbf{9}  & \textbf{VB}  & 4   & V     & 5.9   & 15 & Base motora elementar; transferência de peso sobre pés grandes sem força isolada de dedos. \\
\textbf{11} & \textbf{V0-} & 4+  & VI    & 5.10b & 17 & Consciência de quadril próximo à parede; bordas positivas largas e pés direcionais. \\
\rowcolor{rowgray}
\textbf{12} & \textbf{V0}  & 5   & VI    & 5.10c & 18 & Equilíbrio espacial; agarras oblíquas positivas, oposição e trocas de pé precisas. \\
\textbf{13} & \textbf{V0+} & 5+  & VIsup & 5.10d & 19 & Transição intermediária; negativos leves ou placas técnicas com apoios de pé reduzidos. \\
\rowcolor{rowgray}
\textbf{14} & \textbf{V1}  & 6A  & VIIa  & 5.11a & 20 & Força isométrica nos flexores dos dedos; bordas retas, pinças médias e leitura de sequência. \\
\textbf{15} & \textbf{V2}  & 6A+ & VIIb  & 5.11b & 21 & Ganchos de calcanhar (\textit{heel hook}), bandeiras (\textit{flagging}) e pequenos botes. \\
\rowcolor{rowgray}
\textbf{16} & \textbf{V3}  & 6B  & VIIc  & 5.11c & 22 & Negativos moderados e placas exigentes; cruzamentos sob tensão, abaulados e compressão. \\
\rowcolor{rowgray}
\textbf{17} & \textbf{V3+} & 6B+ & VIIc  & 5.11d & 23 & Sequências contínuas de 5 a 8 movimentos intensos sem descanso intermediário. \\
\rowcolor{rowgray}
\textbf{18} & \textbf{V4}  & 6C  & VIIIa & 5.12a & 24 & Nível avançado inicial; regletes de 15--20\,mm, tensão de tronco e bloqueios altos. \\
\textbf{19} & \textbf{V5}  & 6C+ & VIIIb & 5.12b & 25 & Bloqueios unilaterais profundos, \textit{toe hooks} em tetos e dinâmicos longos. \\
\rowcolor{rowgray}
\textbf{20} & \textbf{V6}  & 7A  & VIIIc & 5.12c & 26 & Alto desempenho; coordenação dinâmica complexa, regletes vivos e compressão em volumes. \\
\textbf{21} & \textbf{V7}  & 7A+ & IXa   & 5.12d & 27 & Nível expert; regletes mínimos ($\le 12$\,mm) e alta taxa de produção de força (RFD). \\
\rowcolor{rowgray}
\textbf{22} & \textbf{V8}  & 7B  & IXb   & 5.13a & 28 & Alta performance competitiva; força digital extrema, mobilidade de quadril e precisão. \\
\rowcolor{rowgray}
\textbf{24} & \textbf{V9}  & 7C  & IXc   & 5.13c & 29 & Entrada na elite esportiva; alta intensidade sustentada sem margem para ajuste postural. \\
\rowcolor{rowgray}
\textbf{26} & \textbf{V10} & 7C+ & Xb    & 5.14a & 31 & Padrão de finais nacionais/internacionais e projetos clássicos em rocha natural. \\
\bottomrule
\end{tabularx}
\end{table}

\clearpage
\subsection*{3.1 Matriz de Referência de Circuitos (Nove Faixas Cromáticas)}

% SOURCE [INTERPRETATION]: Síntese aplicada das equivalências IRCRA (Draper et al., 2016) e CBME DT-2016/01 com divisão da faixa de entrada em Branco-Vazado e Branco.
\begin{table}[!htbp]
\centering
\caption{Matriz de Referência de Circuitos por Cores (Nove Faixas) e Equivalência Multissistêmica.}
\label{tab:circuitos_referencia}
\scriptsize
\renewcommand{\arraystretch}{1.15}
\begin{tabularx}{\linewidth}{p{2.55cm}ccccp{1.8cm}X}
\toprule
\textbf{Circuito / Cor} & \textbf{V-Scale} & \textbf{FB} & \textbf{IRCRA} & \textbf{CBME} & \textbf{RIC} & \textbf{Padrão Motor e Referência Prática} \\
\midrule
\colorbox{tagwhitehollow}{\strut\color{black}\textbf{0. Branco-Vazado}} & \textbf{VB / V0} & 4--5 & 8--12 & IV--VI & \texttt{R1 I1 C1} & \textbf{Iniciação Absoluta}: Agarras anatômicas francas (\textit{jugs}), movimento tipo escada em parede positiva ou vertical. \\
\rowcolor{rowgray}
\colorbox{tagwhite}{\strut\color{black}\textbf{1. Branco}} & \textbf{V0 / V1} & 5--6A & 12--14 & VI--VIIa & \texttt{R1 I1-2 C1-2} & \textbf{Transição Básica}: Introdução a bordas positivas, equilíbrio espacial e primeiros passos de oposição. \\
\colorbox{tagyellow}{\strut\color{black}\textbf{2. Amarela}} & \textbf{V1 / V2} & 6A--6A+ & 14--15 & VIIa/b & \texttt{R1-2 I2 C2} & \textbf{Iniciante Consolidado}: Centro de gravidade, pés precisos, primeiros \textit{heel hooks} e posicionamento lateral. \\
\rowcolor{rowgray}
\colorbox{taggreen}{\strut\color{white}\textbf{3. Verde}} & \textbf{V2 / V3} & 6A+--6B+ & 15--17 & VIIb/c & \texttt{R2 I2-3 C2-3} & \textbf{Intermediário}: Negativos moderados, placas em ponta de pé, pinças e cruzamentos controlados. \\
\colorbox{tagblue}{\strut\color{white}\textbf{4. Azul}} & \textbf{V3 / V4} & 6B--6C & 16--18 & VIIc/VIIIa & \texttt{R2-3 I3 C3} & \textbf{Intermediário-Avançado}: Compressão em volumes, regletes médios (18--22\,mm) e tensão abdominal. \\
\rowcolor{rowgray}
\colorbox{tagred}{\strut\color{white}\textbf{5. Vermelha}} & \textbf{V4 / V5} & 6C--6C+ & 18--19 & VIIIa/b & \texttt{R3 I3-4 C3-4} & \textbf{Avançado}: Bloqueios pronunciados (\textit{lock-offs}), saídas de teto e botes de coordenação. \\
\colorbox{tagblack}{\strut\color{white}\textbf{6. Preta}} & \textbf{V5 / V6} & 6C+--7A & 19--20 & VIIIb/c & \texttt{R3-4 I4 C4} & \textbf{Avançado Consolidado}: Patamar de Font 7A; regletes incisivos (12--15\,mm) e compressão ampla. \\
\rowcolor{rowgray}
\colorbox{tagsilver}{\strut\color{black}\textbf{7. Prata}} & \textbf{V6 / V7} & 7A--7A+ & 20--21 & VIIIc/IXa & \texttt{R4 I4-5 C4} & \textbf{Alto Rendimento}: Baixa margem de erro, força isométrica elevada e coordenação tridimensional. \\
\colorbox{taggold}{\strut\color{black}\textbf{8. Dourada}} & \textbf{V7 / V8+} & $\ge$ 7A+/7B & $\ge$ 21--22 & $\ge$ IXa/b & \texttt{R4-5 I5 C5} & \textbf{Elite Nacional}: Limite biomecânico do ginásio; regletes milimétricos e saltos duplos de competição. \\
\bottomrule
\end{tabularx}
\end{table}

\noindent\textit{Síntese:} A divisão da entrada em \textbf{Branco-Vazado} (VB/V0) e \textbf{Branco} (V0/V1) elimina o salto de cinco degraus IRCRA em uma única cor, preservando a sobreposição gradual a partir do circuito Amarelo.
